% ===========================================================================
% SECTION 1 — INTRODUCTION
% Target: 800-1000 words
% ===========================================================================

\section{Introduction}
\label{sec:introduction}

Wind energy is one of the fastest-growing renewable energy sources, and
the continued growth of installed capacity places increasing demands on
turbine control systems \citep{pao2009}. Above the rated wind speed
(Region~III), a variable-speed, variable-pitch wind turbine must regulate
rotor speed to its rated value by actuating blade pitch, while the
generator torque is held at (or switched into) its rated value
\citep{bianchi2006}. This regulation problem is nonlinear, subject to
actuator rate and range constraints, and driven by a highly stochastic
disturbance (turbulent wind); Model Predictive Control (MPC) is
well-suited to it precisely because it handles input and state
constraints explicitly and can anticipate disturbance effects over a
prediction horizon \citep{rawlings2017}, in contrast to classical
gain-scheduled PI pitch
controllers \citep{bianchi2006,pao2009,laks2009}. The practical obstacle to
nonlinear MPC (NMPC) for this problem is computational: at every control
step, the solver must repeatedly evaluate the internal plant model — and,
for gradient-based solvers, its Jacobian — which is expensive when that
model is the full nonlinear aerodynamic and structural dynamics of the
turbine \citep{schlipf2013,evans2022}.

\subsection{ML-MPC: State of the Art}
\label{subsec:soa}

A natural response to this computational bottleneck is to replace the
physical plant model inside the MPC prediction step with a learned
surrogate that is cheaper to evaluate. \citet{klein2024,klein2025} recently
demonstrated this approach for wind turbine load-relevant control using a
Local Linear Neuro-Fuzzy Model (LLNFM) surrogate inside an adaptive MPC
framework, validated on a full-scale turbine — the direct reference and
starting point for the two-state modeling and LLNFM implementation adopted
in this work (Section~\ref{sec:system_modeling}). Beyond neuro-fuzzy
surrogates, Gaussian Process (GP) regression has an established track
record as an MPC-internal model precisely because it returns a calibrated
predictive uncertainty alongside its point prediction, enabling cautious
or uncertainty-aware control formulations \citep{kocijan2004,hewing2020} —
a property we exploit directly in the GP-based cost of
Section~\ref{subsec:gp_uncertainty_cost}. Within wind energy specifically,
\citet{liu_wind_2018} use a Mat\'ern-kernel GP to predict short-term wind
speed for preview-based load mitigation, but --- unlike the present
work --- as an external disturbance forecast feeding the controller
rather than as the MPC's internal dynamics model. Data-driven
reinforcement learning has separately been applied to wind turbine
torque/pitch control \citep{xie_data-driven_2023}, offering a
model-free alternative to the surrogate-based MPC approach studied
here; we return to this comparison in
Section~\ref{sec:discussion}. Physics-Informed Neural Networks
(PINNs), which embed a governing differential equation as a soft
constraint on the training loss, have become a general framework for
combining data-driven flexibility with physical consistency
\citep{raissi2019,karniadakis2021}, and Temporal Convolutional Networks
(TCNs) have been shown to be a competitive, often superior alternative to
recurrent architectures for generic sequence modeling tasks
\citep{bai2018}. Each of these four paradigms — neuro-fuzzy, GP, PINN,
and convolutional/recurrent sequence models — has therefore been
independently validated as a plausible MPC surrogate architecture in some
application domain. Beyond ML surrogates specifically, the broader
engineering practice of MPC continues to evolve rapidly
\citep{schwenzer2021}, and linear-quadratic-optimal alternatives to
nonlinear MPC for wind turbine power tracking, without a learned
internal model, have also recently been proposed and validated in
OpenFAST \citep{grapentin2022,grapentin2024} --- a useful point of
comparison for the computational-cost trade-offs discussed in
Section~\ref{sec:mpc_framework}, though outside this paper's
ML-surrogate scope.

What is missing, to our knowledge, is a systematic comparison of these
paradigms \emph{within a single closed-loop MPC framework, on the same
plant, under the same solver}, paired with an explicit operational
feasibility analysis of running each surrogate inside a real-time SQP
solve rather than only an open-loop accuracy comparison. Prior work
typically evaluates one surrogate family in isolation
\citep{klein2024,kocijan2004,hewing2020}; the interaction between
surrogate architecture and the specific numerical behavior of a
gradient-based NMPC solver — which, as we show in
Section~\ref{subsec:operational_feasibility}, can render an
otherwise-accurate surrogate practically unusable — is, to our knowledge,
not addressed in the existing wind turbine ML-MPC literature.

\subsection{Contributions}
\label{subsec:contributions}

This paper makes four contributions:

\begin{enumerate}
\item A systematic comparison of six ML surrogate architectures (LLNFM,
LSTM, GP, PINN, TCN, SW-MLP) as internal state predictors for nonlinear
MPC of wind turbine rotor speed, evaluated both open-loop (one-step
prediction accuracy) and closed-loop (tracking performance under
turbulent wind), across two model generations (V1, V2).
\item An operational feasibility analysis documenting a hard
incompatibility between \texttt{dlarray}-based deep-learning surrogates
and SQP-based \texttt{nlmpc} in MATLAB R2024a, driven by non-deterministic
JIT recompilation rather than a soft accuracy-cost trade-off, with
per-call and per-solve timing evidence.
\item A Region~II/III generator torque switching correction for the
two-state plant model, calibrated at the turbine's actual rated operating
point rather than its theoretical aerodynamic optimum, without which
long-horizon closed-loop simulation enters an irrecoverable operating
trap — a modeling detail we show is not a minor implementation nicety but
a precondition for valid training data and closed-loop evaluation.
\item A Gaussian-Process-based uncertainty-penalized MPC cost, made
statistically meaningful through Matérn~5/2 kernel selection and Bayesian
hyperparameter optimization, extending the cautious-MPC literature
\citep{hewing2020,kocijan2004} to Region~III wind turbine speed
regulation.
\end{enumerate}

\subsection{Paper Organization}
\label{subsec:organization}

Section~\ref{sec:system_modeling} presents the plant, torque-switching,
and wind models. Section~\ref{sec:dataset_generation} describes the
training data generation protocol. Section~\ref{sec:ml_surrogates}
details all six surrogate architectures and their open-loop accuracy.
Section~\ref{sec:mpc_framework} presents the NMPC formulation, surrogate
integration, and operational feasibility analysis. Section~\ref{sec:closed_loop_results}
reports closed-loop V1 and V2 results, their comparison, and a Monte Carlo
robustness analysis. Section~\ref{sec:discussion} discusses the main
findings, limitations, and directions for future work, and
Section~\ref{sec:conclusion} concludes.
